% Bagean: metodhe
% Bab: induksi
% Pérangan: induksi-kuwat

\documentclass[../../../include/open-logic-section]{subfiles}

\begin{document}

\olfileid{mth}{ind}{str}

\olsection{Induksi Kuwat}

Ing prinsip induksi kang wis dirembug,
kita mbuktekaké $P(0)$ lan uga yèn
$P(k)$ mula $P(k+1)$. Ing pérangan
kapindho, kita nganggep $P(k)$ bener
lan migunakaké anggepan iki kanggo
mbuktekaké $P(k+1)$. Kanthi cara kang
padha, kita bisa nganggep $P(k-1)$
lan migunakaké kanggo mbuktekaké
$P(k)$. Kang wigati yaiku kita bisa
nalar saka sembarang bilangan menyang
bilangan panerusé, yaiku mbuktekaké
pratelan kanggo sembarang bilangan
kanthi nganggep pratelan iku bener
kanggo bilangan sadurungé.

Ana varian prinsip induksi kang
ora mung nganggep pratelan bener
kanggo bilangan sadurungé $k-1$
saka $k$, nanging kanggo kabèh
bilangan kang luwih cilik tinimbang~$k$,
banjur migunakaké anggepan iki
kanggo mbuktekaké pratelan tumrap~$k$.
Iki uga mènèhi $P(n)$ kanggo
saben~$n \in
\Nat$. Sawisé $P(0)$ kabuktèkaké,
kita wis nuduhaké yèn $P$ bener
kanggo kabèh bilangan kang luwih
cilik tinimbang~$1$. Yèn kita
ngerti yèn saka $P(l)$ kanggo
saben $l<k$ katut $P(k)$,
iki mligi bener kanggo $k=1$.
Dadi kita nyimpulaké $P(1)$.
Saiki $P(0)$ lan $P(1)$ wis
kabuktèkaké, yaiku $P(l)$ kanggo
saben $l<2$. Amarga kita uga
nduwèni kondisional yèn $P(l)$
kanggo saben $l<2$ mula $P(2)$,
kita bisa nyimpulaké~$P(2)$,
lan saterusé.

Sajatiné, yèn kita bisa mbuktekaké
kondisional umum ``kanggo saben $k$,
yèn $P(l)$ kanggo saben $l<k$
mula $P(k)$'', kita ora perlu
mbuktekaké $P(0)$ kanthi kapisah,
amarga pratelan iki wis katut.
Élinga yèn pratelan umum kaya
``kanggo saben $l<k$, $P(l)$''
bener yèn ora ana $l<k$.
Iki kasus kuantifikasi kang bener kanthi vakum:
``kabèh $A$s iku $B$s'' bener
yèn ora ana $A$s, lan
$\lforall[x][(!A(x) \lif !B(x))]$
bener yèn ora ana $x$ kang
nyukupi~$!A(x)$. Ing kasus iki,
wujud formalé yaiku
``$\lforall[l][(l < k \lif P(l))]$'',
lan iki bener yèn ora ana
$l < k$. Nalika $k=0$, kahanan
iki kelakon: ora ana $l<0$,
mula ``kanggo saben $l<0$,
$P(l)$'' bener, apa waé $P$.
Dadi bukti kanggo ``yèn $P(l)$
kanggo saben $l<k$ mula $P(k)$''
kanthi otomatis mbuktekaké~$P(0)$.

Varian iki migunani yèn pratelan
kanggo~$k$ ora bisa dibuktèkaké
mung saka pratelan kanggo $k-1$,
nanging mbutuhaké anggepan yèn
pratelan mau bener kanggo siji
utawa luwih $l<k$.

\end{document}
